% Experiment Details appendix. Requires amsmath, amssymb, graphicx, booktabs.
% Include after the parent manuscript's \appendix.
% Override \appendixfigdir for the parent manuscript's asset directory.
\providecommand{\appendixfigdir}{paper_plotting/paper_figures/appendix}

\section{Experiment Details}
\label{app:experiment-details}

This appendix describes the data partitions, training procedures, estimators,
and numerical optimization used in the experiments. Section~\ref{app:celeba-training-details} gives the details for the CelebA experiment, Section~\ref{app:dsprites-training-details} gives the details for the dSprites experiment, and Section~\ref{app:gaussian-solver} describes the algorithm used for the Gaussian covariance optimization \eqref{eq:program}. The main crossover and recovery figures remain in the main text.

\paragraph{Scope and model selection.}
We describe the frozen-backbone CelebA experiment and the main dSprites
experiment used in the paper figures. Within each experiment, pretexts share
the representation architecture, optimization settings, data-partition
protocol, and rate or regularization grid. Test losses are used for reporting,
not for the checkpoint or hyperparameter selection procedures described
below. Architectural dimensions, learning rates, epoch budgets, and numerical
regularization constants are fixed configuration settings. The recorded
experiments do not contain a broader search establishing their optimality.
We distinguish these fixed settings from the explicitly implemented
validation-based choices.

\paragraph{Computational resources.}
The experiments were conducted on computing infrastructure equipped with
four NVIDIA H100 GPUs.

\subsection{CelebA experiment}
\label{app:celeba-training-details}

\paragraph{Data and frozen features.}
We sample without replacement 10,000 training images, 2,000 validation images,
and 2,000 test images from the official CelebA partitions of sizes
162,770, 19,867, and 19,962, respectively. Subsampling uses NumPy generators
with seeds 42, 43, and 44 for the three partitions. The validation subset
is randomly divided, using seed $42+9182$, into 1,000 images for encoder
checkpoint selection and 1,000 images for selection among pretexts. The
training and test subsets remain unchanged across the five representation
seeds $\{42,43,44,45,46\}$. These seeds therefore do not constitute five
independent data splits.

An ImageNet-1K V1 pretrained ResNet-18, with its final classifier removed,
provides 512-dimensional features. The backbone is frozen and evaluated in
inference mode. Anchor images use deterministic bilinear resizing to a shorter edge of 256,
a $224\times224$ center crop, and channel normalization with means
$(0.485,0.456,0.406)$ and standard deviations $(0.229,0.224,0.225)$. PCA and whitening are fitted exclusively to training
anchor features and retain at most 32 eigen-directions with eigenvalues
exceeding $\max(10^{-5}\lambda_{\max},10^{-12})$. The fitted transform is
applied unchanged to validation and test images. Target-view features use
separately fitted training-only PCA/whitening transforms, also of dimension
at most 32. Feature extraction uses batches of 64 and two data-loader workers.

\paragraph{Tasks and pretexts.}
The four equally weighted evaluation attributes are \texttt{Smiling},
\texttt{Eyeglasses}, \texttt{Blond\_Hair}, and \texttt{Young}. The specialist
shown in the current figures predicts five disjoint source attributes:
\texttt{Mouth\_Slightly\_Open}, \texttt{High\_Cheekbones},
\texttt{Wavy\_Hair}, \texttt{Straight\_Hair}, and
\texttt{Receding\_Hairline}. The generalist reconstructs the anchor feature
vector, rather than image pixels. The complete eight-pretext bank additionally
contains lower-region prediction, grayscale-view prediction, separate
expression and hair attribute objectives, an equally weighted
semantic/reconstruction objective, and equally weighted upper/lower-region
prediction. The task-aware baseline uses the four evaluation attributes
on the training partition. It is distinct from the source-label pretexts.

Two reproducible target views per configured image policy are cached.
Target augmentations use a random resized crop to $224\times224$ with crop
area fraction in $[0.9,1]$ and horizontal-flip probability $0.5$.
Grayscale and region-prediction views use brightness, contrast, and saturation
jitter of $0.1$ and hue jitter of $0.025$. Upper and lower crops retain
horizontal fractions $[0.08,0.92]$ and vertical fractions $[0.05,0.62]$ and
$[0.40,0.97]$, respectively, with a vertical shift sampled from
$[-0.025,0.025]$. Views use ImageNet channel normalization. The cached views
are reused during bottleneck training, rather than generated afresh each epoch.

\paragraph{Gaussian surrogates and channels.}
For each scalar attribute, we fit a ridge regression with an intercept and
ridge coefficient $10^{-3}$ on whitened training features. Five-fold shuffled
cross-fitting, with seed 42, estimates the residual mean squared error.
The final mean predictor is refitted on all training examples. With fitted
slope $b_t$ and residual variance
$\hat\sigma_t^2=\max(\mathrm{MSE}_{\mathrm{OOF},t},10^{-4})$, the task matrix is
$H_t=b_tb_t^\top/\hat\sigma_t^2$. The cross-fitting folds share the unlabeled
PCA/whitening transform fitted on the complete training partition. Preprocessing
is not refitted separately within each regression fold. Feature-prediction
surrogates are obtained from training cross-covariances and canonical
correlations, clipped at $0.99$ before conversion to Gaussian signal-to-noise
ratios. Source-attribute surrogates use the same ridge/cross-fitting procedure.

The nominal Gaussian rate grid is
$\{0,0.25,1,2,4,8,16,32\}$ bits. Given residual covariance $K$, the channel
is $C=AX+\varepsilon$, with independent standard Gaussian noise and
$A^\top A=K^{-1}-I$. Pure-noise output directions are removed. Zero rate
is represented by a constant feature. Two independent channel-noise draws
per image are used for probe fitting and scoring. Random-number seeds are
shared across candidates and budgets. The nominal rate is mutual information
for the fitted whitened Gaussian model, not an empirical measurement of
mutual information for the non-Gaussian image features.

\paragraph{Finite-code training.}
Code budgets are $B\in\{0,2,4,8,16,32\}$ bits per image. For $B>0$, the
encoder is a $32\!\rightarrow\!128\!\rightarrow\!B$ multilayer perceptron
(with the actual retained PCA dimension replacing 32 if necessary), using
GELU after its hidden layer. Forward passes produce hard bits by thresholding
logits at zero. During training, a straight-through estimator uses the
sigmoid derivative while retaining hard bits in the forward pass. Evaluation
uses deterministic hard bits. Each prediction head has architecture
$B\!\rightarrow\!128\!\rightarrow\!m$ with GELU, where $m$ is the target
dimension. Attribute heads minimize mean binary cross-entropy. Feature heads
minimize Gaussian negative log likelihood per target coordinate, with a
learned full residual covariance parameterized by a Cholesky factor whose
diagonal is $\operatorname{softplus}(a)+10^{-3}$. Composite objectives use
separate heads and a shared encoder. At zero bits, no encoder is trained.

Each nonzero-budget encoder is trained independently for 30 epochs using
AdamW, learning rate $10^{-3}$, weight decay $10^{-4}$, batch size 256, and
gradient-norm clipping at 10. Training shuffles images every epoch and samples
one cached target draw per image and update. The checkpoint minimizing the
corresponding hard-code validation objective, averaged over all cached
validation target draws, is retained. This is best-checkpoint selection over
a fixed 30-epoch run, not early termination. The trained task-aware code
baseline uses the same procedure and is not claimed to be globally optimal.

\paragraph{Downstream probes and selection.}
Each attribute is evaluated with an independent logistic regression fitted
on representations of the 10,000 training images. An internal randomized
80/20 training partition selects $C\in\{0,0.1,1,10\}$ by binary log loss.
$C=0$ denotes an intercept-only predictor. The selected model is refitted on
all training representations. Nonconstant logistic models use L-BFGS,
tolerance $10^{-7}$, and at most 300 iterations. For stochastic channels,
regularization selection uses one draw per training image, followed by
refitting on both draws, avoiding a split of replicas of the same image
between fitting and selection. These CelebA probes reuse representation-training
images. Unlike the dSprites experiment, they do not have an independent
probe-training image partition.

The rate-dependent pretext rule minimizes mean downstream loss on the
separate 1,000-image selection subset, averaged across representation seeds.
The fixed-pretext rule minimizes that selection loss averaged across budgets
and seeds. The Gaussian-predicted rule instead minimizes the fitted Gaussian
surrogate distortion. None of these rules uses test losses. The figures'
explicit specialist/generalist comparison is a specified pair, not a
per-budget test-selected comparison. Reported measured excess NLL subtracts
the fitted uncompressed probe's loss, rather than a known Bayes loss.

The decoder-capacity diagnostic replaces logistic probes with one-hidden-layer
64-unit ReLU classifiers. It uses Adam, learning rate $10^{-3}$, L2 parameter
$10^{-4}$, batches of at most 256, at most 100 epochs, and patience 10 on
checkpoint-validation NLL. Validation improvements smaller than $10^{-8}$
are not counted. Hyperparameters are shared across pretexts.

\paragraph{Uncertainty estimates.}
CelebA measured Gaussian-channel bands use 2,000 paired test-image bootstrap
resamples conditional on the fitted seed runs. Learned-code bands additionally
resample training seeds, except in the decoder-capacity diagnostic, whose
bands resample test images conditional on the five fitted seeds.
Repeated channel draws are averaged at the image level before resampling.

A separate bootstrap for the Gaussian proxy tasks requests 200 training-image
resamples and refits PCA/whitening and task/pretext moments in each resample,
grouping bootstrap copies of an original image in the same cross-fitting fold.
This is conditional on the frozen backbone and cached view bank.

\paragraph{Crossover and ranking estimators.}
For a generalist-minus-specialist difference curve on the observed rate grid,
the crossover is the first positive-to-negative zero crossing of its
piecewise-linear interpolant. For bootstrap curves, the median and
2.5/97.5 percentiles are computed only over resamples with a crossing.
The crossing fraction and the number of resamples with multiple crossings
are reported separately. No extrapolation beyond the measured grid is used.
Pretexts are ranked by distortion at each nonzero rate. Endpoint instability
is summarized by tie-corrected Kendall $\tau_b$ and the number of discordant
pairs between the lowest and highest nonzero rates. Zero rate is excluded
because all pretexts tie there.

\paragraph{Diagnostics.}
Figures~\ref{fig:app-achieved-rate} and~\ref{fig:app-whitening} assess code
occupancy and the fitted feature coordinates. The plots use retained run
data. No neural networks were retrained to assemble this appendix.

The pair compared in the main text consists of the five-source-attribute
specialist \texttt{joint\_attributes} and the anchor-feature reconstruction
generalist. Figure~\ref{fig:app-full-bank} reports loss curves for every
configured pretext in both compression mechanisms, complementing the
two-panel rank comparison in Figure~\ref{fig:app-rank-bank}.

Figures~\ref{fig:app-per-task} and~\ref{fig:app-per-task-codes}
resolve the aggregate comparison into individual evaluation attributes,
showing that the crossover need not occur uniformly across tasks.
Figure~\ref{fig:app-subspace-allocation} examines information allocation
within and outside the specialist span.
Figure~\ref{fig:app-proxy-fidelity} assesses the association between
Gaussian surrogate distortion and measured excess NLL.
Figure~\ref{fig:app-decoder-capacity} evaluates sensitivity to replacing
linear probes with MLP probes in both compression mechanisms.
Table~\ref{tab:app-theory-data} reports the fitted theoretical quantities
alongside the observed crossovers.

\begin{figure*}[t]
  \centering
  \includegraphics[width=.46\linewidth]{\appendixfigdir/f2a_code_entropy.pdf}\par\medskip
  \includegraphics[width=.88\linewidth]{\appendixfigdir/f2a_code_usage.pdf}
  \caption{\textbf{Finite-code rate accounting.}
    Top: empirical joint code entropy versus nominal budget for
    300 fitted candidate/budget/seed cells. The identity line
    denotes equality. The maximum observed entropy is 10.96578 bits,
    equal to the finite-test-set ceiling $\log_2(2000)=10.96578$.
    Thus this estimator cannot resolve a 16- or 32-bit population entropy.
    Bottom: test-image codeword-occupancy histograms for the five-attribute
    specialist and reconstruction generalist at all six budgets,
    $B\in\{0,2,4,8,16,32\}$, using seed 42.
    Unobserved codewords are omitted. Sparse occupancy can bias plug-in
    entropy downward. The histograms use logarithmic occupancy and symlog
    count axes to display both singleton and heavily reused codewords.}
  \label{fig:app-achieved-rate}
\end{figure*}

\begin{figure*}[t]
  \centering
  \includegraphics[width=\linewidth]{\appendixfigdir/f2b_covariance_spectrum.pdf}
  \caption{\textbf{Training covariance and whitening.}
    Left: the 512 eigenvalues of the empirical covariance of the 10,000
    training anchor features before PCA, using normalization $1/n$.
    Eigenvalues range from 0.004730 to 16.451719.
    Right: covariance eigenvalues after applying the saved 32-dimensional
    PCA/whitening transform to the same training features. Their maximum
    absolute deviation from one is $5.11\times10^{-15}$, confirming the numerical
    training-coordinate normalization. These are fitted spectra without
    uncertainty bands. No shrinkage-sensitivity sweep is available. This
    figure does not establish robustness to a different shrinkage choice.}
  \label{fig:app-whitening}
\end{figure*}


\begin{figure*}[t]
  \centering
  \includegraphics[width=\linewidth]{\appendixfigdir/f5_full_pretext_bank.pdf}
  \caption{\textbf{Complete CelebA pretext-bank comparison.}
    Both panels show the eight configured pretexts: expression attributes,
    hair attributes, joint attributes, grayscale prediction, lower-region
    prediction, reconstruction, semantic reconstruction, and two-region
    prediction. The task-aware and average-task baselines are also included.
    Curves show five-seed mean task-averaged measured excess NLL relative
    to the uncompressed probe, with markers at each evaluated rate. Bands
    are omitted here for legibility. Colors and markers identify the same
    objectives across panels. The task-aware curve is dotted orange in
    both panels, and the average-task reference is dashed. The left axis
    uses nominal Gaussian rate, the right code budget. The different
    curve shapes reveal objective-dependent changes in relative performance.}
  \label{fig:app-full-bank}
\end{figure*}

\begin{figure*}[t]
  \centering
  \begin{minipage}[t]{.49\linewidth}
    \centering
    \textbf{(a) Gaussian channel}\par\smallskip
    \includegraphics[width=\linewidth]{\appendixfigdir/f5_rank_gaussian.pdf}
  \end{minipage}\hfill
  \begin{minipage}[t]{.49\linewidth}
    \centering
    \textbf{(b) Learned binary codes}\par\smallskip
    \includegraphics[width=\linewidth]{\appendixfigdir/f5_rank_codes.pdf}
  \end{minipage}
  \caption{\textbf{Pretext-bank rank instability across compressors.}
    Left: Gaussian channel. Right: learned binary codes.
    Both panels use the same pretext order, determined by the Gaussian
    ranking at its lowest nonzero rate.
    Darker cells indicate better ranks, and outlines identify column winners
    including ties. Both panels use the same rank scale. Zero rate and the
    task-aware/average-task references are excluded.
    Gaussian endpoints have $\tau_b=0.357143$, with 9 reversed pairs out of
    28 comparable pairs and a maximum rank change of 6 places.
    Code-budget endpoints have $\tau_b=0.071429$, with 13 reversed pairs
    out of 28 comparable pairs and a maximum rank change of 4 places.
    Relative pretext performance therefore changes with rate under both
    compression mechanisms.}
  \label{fig:app-rank-bank}
\end{figure*}

\begin{figure*}[t]
  \centering
  \includegraphics[width=.94\linewidth]{\appendixfigdir/f5a_gaussian_per_task.pdf}
  \caption{\textbf{CelebA per-attribute Gaussian behavior.}
    Nominal rates are $\{0,0.25,1,2,4,8,16,32\}$ bits. Panels show Smiling,
    Eyeglasses, Blond Hair, and Young. The ordinate is measured excess NLL
    relative to the uncompressed linear probe. Bands are 95\% pointwise
    percentile intervals from 2,000 paired test-image resamples conditional
    on five fitted seeds. At rate 0.25 the specialist has lower mean NLL on
    all four tasks. At rate 32 the generalist is better on Eyeglasses,
    Blond Hair, and Young, while the specialist remains better on Smiling:
    the aggregate reversal is not uniform across tasks.}
  \label{fig:app-per-task}
\end{figure*}

\begin{figure*}[t]
  \centering
  \includegraphics[width=.94\linewidth]{\appendixfigdir/f5a_codes_per_task.pdf}
  \caption{\textbf{CelebA per-attribute learned-code behavior.}
    Code budgets are $B\in\{0,2,4,8,16,32\}$ bits/image. The same four
    evaluation tasks, specialist, generalist, and uncompressed-probe baseline
    are used as in Figure~\ref{fig:app-per-task}. Bands are 95\% pointwise
    percentile intervals from 2,000 paired training-seed/test-image resamples
    using five fitted seeds. These panels expose task-level variation behind
    the aggregate code-budget crossover reported in Table~\ref{tab:app-theory-data}.}
  \label{fig:app-per-task-codes}
\end{figure*}

\begin{figure*}[t]
  \centering
  \includegraphics[width=\linewidth]{\appendixfigdir/f5b_subspace_allocation.pdf}
  \caption{\textbf{Gaussian subspace allocation.}
    The specialist span has dimension 5 within the 32 retained coordinates.
    Inside marginal and outside conditional information add to total
    information by the chain rule in the fitted model. At total rate 32,
    the specialist allocates approximately 32 bits inside and zero outside.
    Reconstruction allocates 5 inside and 27 outside, and the task-aware
    solution approximately 3.103 inside and 28.897 outside.
    Outside marginal information would not generally give this additive
    decomposition for channels coupling the two subspaces.}
  \label{fig:app-subspace-allocation}
\end{figure*}

\begin{figure}[t]
  \centering
  \includegraphics[width=\linewidth]{\appendixfigdir/f5c_proxy_fidelity.pdf}
  \caption{\textbf{Gaussian surrogate fidelity.}
    Ordinary least squares of measured excess NLL on predicted Gaussian
    distortion, with a fitted intercept, gives $R^2=0.9815$ across 80
    objective/rate pairs. The fitted line is
    $\widehat D_{\mathrm{measured}}=0.701817 D_{\mathrm{proxy}}-0.004594$.
    The black line shows this fit. The dashed gray identity line remains
    a calibration reference. The score is descriptive and in-sample, not a
    held-out regression score. Strong linear association coexists with a
    non-unit slope, so it does not imply identity calibration.}
  \label{fig:app-proxy-fidelity}
\end{figure}

\begin{figure*}[t]
  \centering
  \includegraphics[width=\linewidth]{\appendixfigdir/f5d_decoder_capacity.pdf}
  \caption{\textbf{Decoder-capacity sensitivity.}
    Left: Gaussian channel. Right: learned binary codes. Linear and
    64-hidden-unit MLP probes evaluate the same representations in each
    mechanism. The ordinate is generalist minus specialist excess NLL.
    The zero line denotes equal performance. Bands in both panels use
    2,000 paired test-image resamples conditional on the five fitted seeds,
    rather than additional seed resampling. Gaussian crossings are 16.8182
    and 19.2453 nominal bits for linear and MLP probes, respectively.
    The code-panel crossings are 21.2035 and 19.9533 bits/image.
    The two rate axes represent different budgets and are not shared.}
  \label{fig:app-decoder-capacity}
\end{figure*}


\begin{table*}[t]
  \centering
  \input{\appendixfigdir/f5f_theory_to_data.tex}
  \caption{\textbf{Fitted theory-to-data quantities.}
    Values concern the five-source-attribute specialist and reconstruction
    generalist. The first four rows are per-task largest eigenvalues.
    $\Delta=0.0181922$ and $D_{\mathrm{miss}}=0.215195$ are fitted proxy
    quantities. The 16.2447-bit threshold is a sufficient isotropic-reference
    bound, not an equality prediction for the empirical crossover.
    The fitted specialist is anisotropic on its span. The saved strict
    isotropy diagnostic also rejects exact full-space isotropy for the
    reconstruction components. Consequently, the reference threshold is
    not asserted as a certificate for these fitted components.
    Crossover intervals condition on a first positive-to-negative crossing.
    All 2,000 resamples cross in each mechanism, with no multiple crossings.}
  \label{tab:app-theory-data}
\end{table*}


\subsection{dSprites experiment}
\label{app:dsprites-training-details}
\paragraph{Data partitions.}
The primary dSprites experiment uses binary $64\times64$ images from the
official archive, with no resizing or image augmentation. Exact duplicate
pixel arrays are grouped before splitting. Groups are shuffled using
seed $100000+10s$ for experimental seed $s\in\{0,1,2\}$ and assigned intact
to five disjoint folds. Requested sizes are 32,768 representation-training
images, 4,096 representation-validation images, 16,384 probe-training images,
4,096 probe-validation images, and 8,192 test images. Keeping duplicate groups
intact can slightly exceed the requested sizes. In the retained main run,
seed 0 has 4,097 probe-validation images. All other counts equal their
requested values. Splits are identical across pretexts and $\beta$ values
within a seed, but different seeds may reuse images.

The targets are horizontal and vertical position indices, each standardized
using the known uniform 32-level mean $15.5$ and standard deviation
$\sqrt{(32^2-1)/12}$. Normalization is not estimated from held-out labels.
Other rendering factors remain nuisance variation. For $M=4$, task weights
are $(3/4,1/4)$ and independent Gaussian task-noise variances are
$(M^{-2},M^{-4})=(1/16,1/256)$. Task noise is integrated analytically in the
loss, rather than repeatedly sampled as noisy target labels.

\paragraph{Architecture and optimization.}
The image encoder uses three convolutional layers with channels
$1\!\rightarrow\!16\!\rightarrow\!32\!\rightarrow\!32$, kernel sizes
$(5,3,3)$, stride 2, and padding $(2,1,1)$, each followed by SiLU.
The resulting $32\times8\times8$ tensor is flattened and mapped to 128
SiLU units. A linear layer produces the mean and log variance of an
8-dimensional diagonal Gaussian representation. Log variance is clamped
to $[-24,12]$. Representations are sampled by reparameterization.
The task decoder is an $8\!\rightarrow\!128\!\rightarrow\!128\!\rightarrow\!4$
MLP with SiLU hidden activations. It predicts two means and two positive
residual variances, the latter parameterized by softplus plus $10^{-7}$.
The reconstruction decoder is an
$8\!\rightarrow\!128\!\rightarrow\!4096$ MLP with a SiLU hidden layer.

All four pretexts (two specialists, weighted task mixture, and image
reconstruction) minimize
\[
 \mathcal L= D_{\mathrm{pretext}}+
 \beta\,\mathbb E_U\operatorname{KL}_2(q(C\mid U)\Vert\mathcal N(0,I)).
\]
The task loss for target $z_i$, predicted mean $\mu_i$, and total predicted
variance $v_i=\sigma_i^2+v_{\mathrm{res},i}$ is
\[
 \ell_i=\frac{\log(v_i/\sigma_i^2)
       +((z_i-\mu_i)^2+\sigma_i^2)/v_i-1}{2\log 2}.
\]
Specialists minimize their respective task loss. The mixture minimizes
$\tfrac34\ell_1+\tfrac14\ell_2$. Reconstruction minimizes mean Bernoulli
cross-entropy per pixel, converted to bits. Thus reconstruction loss is
normalized by the number of pixels. It is not a summed image log loss.

Each representation is trained from scratch for 100 epochs with Adam,
learning rate $10^{-3}$, batch size 256, no weight decay, and gradient-norm
clipping at 10. The best checkpoint minimizes its own validation
$D_{\mathrm{pretext}}+\beta\mathrm{KL}_2$ objective. Common Monte Carlo
validation draws are reused across epochs by resetting the validation RNG
while restoring the training RNG state afterwards. The regularization grid is
\[
 \beta\in\{0.8,0.4,0.3,0.2,0.15,0.1,0.05,0.02,0.01,
                 0.005,0.002,0.001,0.0005,0.0002\}.
\]
Four pretexts, fourteen regularization values, and three seeds give 168
independent training jobs. The saved learning-rate grid contains only
$10^{-3}$: there is no learning-rate search in this run. The aggregation
code can select among additional learning-rate candidates using the
corresponding pretext validation objective, if such candidates are supplied.
The broad $\beta$ grid samples achieved rates on both sides of the potential
reversal. The retained records do not document a separate optimization
study for this grid. Curves are plotted against measured KL, not $\beta$.

\paragraph{Probes and measurement.}
After representation checkpoint selection, the encoder is frozen and a fresh
MLP with the same architecture as the task decoder is trained on the separate
probe-training fold. Probe training uses equal task weights, Adam at
$10^{-3}$, batch size 256, gradient clipping at 10, and 80 epochs. Equal
weighting prevents the probe objective itself from downweighting the rarer
task. The checkpoint with minimum probe-validation NLL is retained.
Probe validation uses eight fixed channel draws per image. Probe training
uses fresh channel samples. The probe and test RNG seeds are $s+30000$
and $s+40000$, respectively.

Test metrics average eight stochastic-channel losses per image, not predictions before evaluating a loss. We report the analytic test-average KL in bits, prevalence-weighted excess NLL, per-factor MSE and $R^2$, and
$\mathcal I_i=-\tfrac12\log_2\mathrm{MSE}_i$.
The KL is a variational upper bound on mutual information.
The recoverability index is a diagnostic transformation of MSE, not an
allocated number of bits for discrete dSprites factors. Its equality to
$-\tfrac12\log_2(1-R_i^2)$ requires unit empirical target variance.
Per-run standard errors cluster duplicate-image groups. Main paper figures
show every seed and across-seed min--max ranges at fixed $\beta$, rather
than interpreting three-seed ranges as confidence intervals.

Table~\ref{tab:dsprites} summarizes mixture-encoder recovery at selected
settings around the reversal.
Figure~\ref{fig:app-dsprites-accounting} compares summed recoverability
indices with the variational rate bound. These indices do not constitute
an additive information decomposition.
Figure~\ref{fig:app-dsprites-context} compares recovery under both
specialists, reconstruction, and the mixture over the full
regularization grid.

\paragraph{Diagnostics and controls.}
The primary run contains 168 completed jobs, and the role-swap control run contains 42 completed mixture jobs, with 14 regularization values and seeds $\{0,1,2\}$ in each comparison. The role swap retains 100 representation-training epochs, 80 probe-training epochs, learning rate $10^{-3}$, batch size 256, and an 8-dimensional code.

\paragraph{Role-swap control and crossover audit.}
The role swap exchanges the ordered factors to vertical then horizontal,
retaining task-index weights $(3/4,1/4)$ and noise variances $(1/16,1/256)$.
Figure~\ref{fig:app-role-swap} shows that the recovery ordering reverses
when horizontal and vertical positions exchange task roles.
Crossovers are obtained from the piecewise-linear interpolant on the
observed rate grid, without extrapolation. The primary gap uses the first
positive-to-negative crossing. The swapped plot keeps the physical-factor
gap horizontal minus vertical and uses its first negative-to-positive crossing.
The primary and swapped mean-curve crossings are 5.2638 and
5.1938 bits/image, respectively. All three seeds reverse in each
run. Each seed has exactly one expected-direction crossing.
Table~\ref{tab:app-dsprites-crossovers} gives per-seed estimates and the
adjacent observed-rate brackets used for interpolation. The brackets are
not confidence intervals. The Gaussian surrogate reference is
$R^\star=5.2996$ bits.

\begin{table*}
\centering
    \begin{tabular}{lllll}
        \toprule
        $\mathcal R_{\mathrm{VIB}}$ & $\mathcal I_{\mathrm{common}}$ & $\mathcal I_{\mathrm{precise}}$ & Gap & $\sum_i\mathcal I_i$ \\
        \midrule
        1.75 & 1.85 & 0.00 & 1.84 & 1.85 \\
        2.16 & 2.27 & 0.00 & 2.27 & 2.27 \\
        5.10 & 2.71 & 2.62 & \textbf{0.09} & 5.34 \\
        6.19 & 2.96 & 3.47 & \textbf{-0.52} & 6.43 \\
        7.15 & 3.29 & 4.03 & -0.75 & 7.32 \\
        8.38 & 3.74 & 4.62 & -0.87 & 8.36 \\
        \bottomrule
    \end{tabular}
    \caption{Factor recovery for mixture-trained encoders, $M = 4$, averaged over 3 seeds.}
    \label{tab:dsprites}
\end{table*}

\begin{table*}[t]
  \centering
  \input{\appendixfigdir/f4_crossovers.tex}
  \caption{\textbf{Per-seed dSprites crossover audit.}
    Rates are measured variational KL bounds in bits/image. The primary
    gap crosses positive-to-negative. The role-swap physical-factor gap
    crosses negative-to-positive. Every seed has one such crossing.
    Brackets are observed neighboring grid rates, not uncertainty intervals.}
  \label{tab:app-dsprites-crossovers}
\end{table*}

\begin{figure}[t]
  \centering
  \includegraphics[width=\linewidth]{\appendixfigdir/f4a_dsprites_accounting.pdf}
  \caption{\textbf{dSprites recoverability accounting.}
    The 42 mixture runs compare the sum of per-run MSE indices with the
    variational KL bound over the full 14-setting grid. Maximum positive
    index-sum excess is 0.2470 bits. The sum is a recovery diagnostic, not
    a quantitative decomposition of the rate bound. The identity line is
    retained as a visual reference. No linear relationship or
    information-theoretic equality is assumed for these discrete factors.}
  \label{fig:app-dsprites-accounting}
\end{figure}

\begin{figure}[t]
  \centering
  \includegraphics[width=\linewidth]{\appendixfigdir/f4b_role_swap.pdf}
  \caption{\textbf{dSprites role-swap control.}
    Vertical position is common/moderate and horizontal position rare/precise.
    Individual points and thin gap curves show three seeds. Central curves
    show means and error bars give seed min--max ranges at fixed $\beta$.
    The horizontal-minus-vertical gap crosses negative-to-positive at
    5.1938 bits/image. All three seeds reverse. The displayed
    view covers 0--10 bits/image. The dotted reference is
    $R^\star=5.2996$ from the Gaussian surrogate, not an exact
    dSprites prediction. The changed physical-factor ordering supports a
    task-role effect rather than an explanation based solely on one fixed
    spatial direction being harder to recover.}
  \label{fig:app-role-swap}
\end{figure}

\begin{figure*}[t]
  \centering
  \includegraphics[width=\linewidth]{\appendixfigdir/f4d_dsprites_context.pdf}
  \caption{\textbf{Recovery under the full dSprites bank.}
    The two panels show horizontal and vertical recoverability versus
    measured KL for both specialists, reconstruction, and mixture.
    Lines connect three-seed means at each of 14 settings. No confidence
    bands are drawn. At each objective's highest mean achieved rate, the
    horizontal specialist has index pair $(5.663,0.007)$, the vertical
    specialist $(0.081,5.697)$, reconstruction $(3.651,3.773)$, and mixture
    $(5.171,5.626)$. These endpoints occur at different achieved rates and
    are not a matched-rate performance comparison. They illustrate that
    specialists primarily recover their own factor, while reconstruction
    and mixture retain both.}
  \label{fig:app-dsprites-context}
\end{figure*}


\subsection{Gaussian covariance optimization}
\label{app:gaussian-solver}
In this subsection, we explain the algorithm used to compute the optimal Gaussian channel in Proposition \ref{prop:gaussian}.

\paragraph{Program and relevant subspace.}
For positive semidefinite task matrices $H_t\in\mathbb R^{d\times d}$ and
nonnegative weights $w_t$ summing to one, the Gaussian surrogate program is
\begin{equation}
 \min_{0\prec K\preceq I}\quad
 D(K)=\frac{1}{2\log 2}\sum_{t=1}^T w_t
          \log\det(I+F_t^\top K F_t),
 \qquad -\frac{\log\det K}{2\log 2}\le R,
 \label{eq:implemented-covariance-program}
\end{equation}
where $H_t=F_tF_t^\top$. The determinant expression is equivalent to
$\log\det(I+H_t^{1/2}KH_t^{1/2})$ but permits low-rank factors.
Here $K$ is the residual covariance of the whitened Gaussian source given
its representation. Scalar regression targets give rank-one $H_t$.

Before the iterative solve, the code computes a basis $U$ for the numerically
relevant span of $\sum_t w_tH_t$, keeping eigenvalues larger than
$\max(10^{-12},10^{-10}\lambda_{\max})$. It solves in this span, of dimension
$q\le d$, using $\widetilde H_t=U^\top H_tU$, and returns
$K=I+U(\widetilde K-I_q)U^\top$. Irrelevant directions receive no rate.
For the scalar-target task-aware and five-attribute specialist objectives,
$q\le4$ and $q\le5$, respectively. With a zero relevant span the code returns
$K=I$ without iteration. The span truncation is numerical. Certificates
refer to the retained, numerically represented program.

\paragraph{Log-covariance coordinates and projected steps.}
Within the relevant span, the solver parameterizes $K=\exp(-S)$.
It optimizes on
\[
 \mathcal S_R=\{S\succeq0:\operatorname{tr}S=2R\log2\}.
\]
The objective is nondecreasing in $K$, so an optimum can exhaust the rate
budget when the relevant span is nonempty. This parameterization enforces
$K\preceq I$ and the rate constraint without a penalty parameter.
It does not assert ordinary Euclidean convexity of $D(\exp(-S))$.
At positive rate, the main experiment initializes
$S_0=(2R\log2/q)I_q$ separately for each solve. Although the low-level solver
accepts a warm start, the main task/pretext wrappers do not supply one.
At zero rate it returns $K=I_q$ directly.

Let $G_K=\nabla_KD(K)$. The implementation evaluates
\[
 G_K=\frac{1}{2\log2}\sum_t w_t
 F_t(I+F_t^\top K F_t)^{-1}F_t^\top.
\]
If $S=V\operatorname{diag}(s)V^\top$, its gradient is obtained using the
Fr\'echet derivative of the matrix exponential:
\[
 \nabla_SD(\exp(-S))=-V\bigl[C\odot(V^\top G_KV)\bigr]V^\top,
 \quad
 C_{ij}=\begin{cases}
 e^{-s_i},&s_i=s_j,\\
 e^{-\min(s_i,s_j)}\dfrac{1-e^{-|s_i-s_j|}}{|s_i-s_j|},&s_i\ne s_j.
 \end{cases}
\]
The code evaluates the divided difference with \texttt{expm1} and uses its
limiting value when $|s_i-s_j|\le10^{-10}$.

Each candidate step is
$S^+=\Pi_{\mathcal S_R}(S-\eta\nabla_SD)$.
Projection diagonalizes the symmetric argument and projects its eigenvalues
onto the nonnegative simplex of mass $2R\log2$ using sorting and a threshold.
The initial step size is 10. Subsequent iterations use the
Barzilai--Borwein proposal
$\eta=\|\Delta S\|_F^2/\langle\Delta S,\Delta G_S\rangle$
when the denominator exceeds $10^{-18}$, clipped to $[10^{-3},10^7]$.
Backtracking halves the step until
\[
 D(e^{-S^+})\le D(e^{-S})+
 \langle\nabla_SD,S^+-S\rangle+
 \frac{\|S^+-S\|_F^2}{2\eta}+10^{-15}.
\]
At most 60 trials are allowed per iteration. If this limit is exhausted,
the implementation still advances to its last candidate. Convergence must
therefore be assessed from the final certificate, not assumed from the line
search alone.

\paragraph{Stopping certificate.}
The stopping rule uses a geodesic Lagrangian residual, rather than objective
change or a fixed iteration count. Write $K=V\operatorname{diag}(k)V^\top$,
$Q=K^{1/2}$, and $M=2\log2\,QG_KQ$.
Eigen-directions with $k_i<1-10^{-7}$ are classified as active. The rate
multiplier $\lambda$ is the nonnegative mean of the active diagonal entries
of $V^\top MV$. If there are no active directions it is the nonnegative
largest eigenvalue of $M$. On the inactive subspace, the code forms the
positive-semidefinite part of
$(\lambda I-(V^\top MV)_{\mathrm{inactive}})/(2\log2)$
and lifts it to a dual matrix $Z\succeq0$, zero on active directions.
Define
\[
 E=\frac{M-\lambda I}{2\log2}+QZQ,
 \qquad r(K)=-\frac{\log\det K}{2\log2}.
\]
The reported residual-based gap bound is
\begin{equation}
 \widehat\epsilon=
 \max\{0,\lambda(R-r(K))+\operatorname{tr}[Z(I-K)]\}
 +2R\log2\sqrt q\,\|E\|_F.
 \label{eq:implemented-gap}
\end{equation}
This is the implementation's a posteriori certificate for the Gaussian
surrogate program. It is not a statistical error bound for learned image
representations. It is computed at iterations $1,11,21,\ldots$ and again
on the returned iterate. The main tolerance is $10^{-5}$ bits of weighted
surrogate distortion, with at most 3,000 outer iterations. Main-run wrappers
reject an unconverged solve. In the retained task-aware run, all seven
positive-rate solves report 11 iterations and gap bounds below $10^{-5}$.
This is an empirical observation, not an iteration-complexity guarantee.

\paragraph{Single-matrix objectives.}
A pretext represented by one matrix bypasses projected iterations.
For $H=V\operatorname{diag}(h)V^\top$, the code computes
$K=V\operatorname{diag}(k)V^\top$ with
$k_i=\min(1,e^a/h_i)$ on positive eigenvalues and $k_i=1$ otherwise.
Brent's method solves the scalar equation
\[
 \frac12\sum_{i:h_i>0}\max\left\{\frac{\log h_i-a}{\log2},0\right\}=R
\]
in log-water-level coordinates, with absolute tolerance $10^{-12}$.
This avoids loss of relative accuracy when the water level is small.
A multivariate attribute pretext with several scalar task matrices uses the
weighted multi-task solver, not this single-matrix shortcut.

\paragraph{Time complexity.}
Let $q$ be the retained dimension, $T$ the number of matrices in one objective,
$r_t=\operatorname{rank}(\widetilde H_t)$ after numerical factorization,
$N$ the number of outer iterations, and $L\le60$ the maximum number of
backtracking trials used in an iteration. Dense symmetric eigendecompositions,
matrix-exponential reconstruction, gradient basis changes, and PSD/simplex
projection cost $O(q^3)$ each. Simplex sorting adds $O(q\log q)$.
With factored tasks, an objective/gradient evaluation costs
\[
 \mathcal C_H=
 O\!\left(\sum_{t=1}^T(q^2r_t+qr_t^2+r_t^3)\right).
\]
When all tasks have rank at most one, the vectorized implementation reduces
this to $O(Tq^2)$. Certificate evaluation adds $O(q^3+\mathcal C_H)$.
Although backtracking requests only an objective value, the current
\texttt{value} routine also computes its gradient, so the same
$\mathcal C_H$ cost applies to those calls. Consequently, one solve has
iteration cost
\[
 O\!\left(N(1+L)(q^3+\mathcal C_H)\right),
\]
or $O(N(1+L)(q^3+Tq^2))$ for rank-one families and
$O(N(1+L)Tq^3)$ in the full-rank case. These are arithmetic-operation bounds
conditional on $N$ and $L$. The implementation establishes no worst-case
bound on $N$ as a function of the requested gap tolerance.

Preprocessing also matters. Forming the weighted matrix and its relevant-span
basis costs $O(Td^2+d^3)$. In the current wrappers, the direct, unoptimized
three-operand \texttt{einsum} for $U^\top H_tU$ has contraction count
$O(Td^2q^2)$. A staged multiplication could reduce this to
$O(T(d^2q+dq^2))$, but that optimization is not implemented here.
Factorizing the reduced matrices costs $O(Tq^3)$. Lifting the solution costs
$O(dq^2+d^2q)$. These setup operations are repeated by the main wrappers
at each rate rather than cached across the grid.
The single-matrix shortcut costs $O(d^3+Jd)$, where $J$ is the number of
scalar root-function evaluations, including dense reconstruction of $K$.

\paragraph{Space complexity and repeated solves.}
The solver retains dense task matrices even for rank-one families. Its
resident matrix storage is therefore $O(Tq^2)$, with factors requiring
$O(q\sum_t r_t)$ and a constant number of $q\times q$ work arrays and
iterates. Including the original matrices and lifted covariance in the
wrappers gives $O(Td^2+Tq^2+dq)$ storage for one solve. It would be incorrect
to describe the present rank-one implementation as using only $O(Tq)$
space. Storing solutions for $G$ rates additionally requires $O(Gd^2)$
per candidate. The solve itself does not depend on the number of training
images once the task matrices have been formed. Fitting features, moments,
and neural representations has separate costs. Rate grids and bootstrap
replications repeat the solve, so their work is additive over all
candidate/rate/resample combinations, including the corresponding refitting
costs for the moment bootstrap.
